%==============================================================================
%  Compatibilidad entre modalidades de resolución colaborativa
%  Borrador de trabajo. Compilar con: pdflatex cpp_framework.tex  (x2)
%  Sin bibliografía por decisión del autor; ver \pendiente en §2.
%==============================================================================
\documentclass[11pt,a4paper]{article}

%------------------------------- codificación --------------------------------
\usepackage[utf8]{inputenc}
\usepackage[T1]{fontenc}
\usepackage[spanish,es-noshorthands,es-noquoting]{babel}
\addto\captionsspanish{\renewcommand{\tablename}{Tabla}}

%------------------------------- matemática ----------------------------------
\usepackage{amsmath,amssymb,amsthm}

%------------------------------- presentación --------------------------------
\usepackage{booktabs}
\usepackage{xcolor}
\usepackage{geometry}
\usepackage[hidelinks]{hyperref}
\geometry{margin=2.6cm}

%------------------------------- entornos ------------------------------------
\theoremstyle{definition}
\newtheorem{definicion}{Definición}[section]
\newtheorem{supuesto}[definicion]{Supuesto}
\newtheorem{requisito}[definicion]{Requisito}
\theoremstyle{plain}
\newtheorem{proposicion}[definicion]{Proposición}
\newtheorem{corolario}[definicion]{Corolario}
\theoremstyle{remark}
\newtheorem{observacion}[definicion]{Observación}

%------------------------------- macros --------------------------------------
\DeclareMathOperator{\Var}{Var}
\DeclareMathOperator{\Cl}{Cl}
\DeclareMathOperator{\clip}{clip}
\newcommand{\ind}[1]{\mathbf{1}\!\left[#1\right]}
\newcommand{\solver}{\mathcal{S}}
\newcommand{\Info}{\mathcal{I}}
\newcommand{\Real}{\mathbb{R}}
\newcommand{\Esp}{\mathbb{E}}
\newcommand{\modAA}{\textsc{aa}}
\newcommand{\modHA}{\textsc{ha}}
\newcommand{\pendiente}[1]{\textcolor{red}{\textbf{[pendiente:} #1\textbf{]}}}

%------------------------------- portada -------------------------------------
\title{\bf Compatibilidad entre modalidades de resolución colaborativa\\
  con partición de información:\\
  \large un marco de medición para el diálogo agente--agente\\
  y el diálogo agente--estudiante}

\author{Luciano Avegno Cepeda\\
  \small [afiliación]\\
  \small \texttt{lucianoavegno@gmail.com}}

\date{Borrador de trabajo --- \today}

\begin{document}
\maketitle

\begin{abstract}
\noindent
Un problema con partición de información reparte entre dos participantes datos
complementarios tales que ninguno alcanza la solución por separado. La misma
instancia puede resolverse en dos modalidades: dos agentes artificiales
dialogando entre sí, o un agente y un estudiante universitario. Este trabajo
sostiene que ambas modalidades son objetos de estudio de pleno derecho y que la
pregunta relevante no es cuál aproxima a cuál, sino \emph{en qué medida una misma
instancia induce la misma estructura colaborativa bajo ambas}. Formalizamos esa
noción como \emph{compatibilidad de modalidad} y la descomponemos en tres
componentes separables por equiparación lineal: sesgo sistemático, distorsión de
escala e incompatibilidad idiosincrásica. Antes de eso reformulamos la batería de
indicadores de colaboratividad, sustituyendo las cantidades probabilísticas no
observables del marco de partida por estimadores empíricos definidos respecto de
un solver de referencia declarado; demostramos en particular que el balance de
carga informativa admite, para dos participantes, una forma cerrada equivalente al
valor de Shapley que sólo requiere cuatro proporciones binomiales y elimina la
entropía condicional del marco. El diseño experimental propuesto es simétrico: un
estudio agente--agente sobre el banco completo y un estudio agente--estudiante con
un brazo extensivo, orientado a la equiparación, y un brazo intensivo,
dimensionado para pruebas de equivalencia estadística. Derivamos la potencia
requerida y mostramos que la equivalencia a nivel de ítem exige aproximadamente
$69$ episodios por modalidad con un margen de media desviación estándar, lo que
delimita con precisión qué afirmaciones el diseño puede sostener y cuáles no.

\vspace{0.6em}
\noindent\textbf{Palabras clave:} resolución colaborativa de problemas,
información asimétrica, interdependencia epistémica, agentes conversacionales,
pruebas de equivalencia, educación matemática.
\end{abstract}

%==============================================================================
\section{Introducción}
%==============================================================================

Considérese un problema de matemática universitaria escrito de modo que su
enunciado público sea insuficiente, y que la información faltante esté repartida
en dos datos privados: uno en poder de cada participante. Ninguno puede resolverlo
solo; juntos determinan la respuesta de forma única. La estructura obliga a
comunicar, y al comunicar obliga a explicitar el propio razonamiento.

Esa misma instancia admite al menos dos modalidades de resolución. En la primera,
dos agentes artificiales reciben cada uno un dato y dialogan entre sí. En la
segunda, un estudiante universitario recibe un dato y un agente artificial recibe
el otro. Las dos producen transcripts, las dos terminan (o no) en una respuesta, y
las dos son mensurables con el mismo instrumento.

La tentación metodológica es tratar una como sustituto de la otra: usar el diálogo
agente--agente, que es barato y paralelizable, para predecir lo que ocurriría con
estudiantes, que es caro y lento. Esa formulación presupone que la modalidad
humana es el criterio y la artificial su aproximación. Nosotros la rechazamos por
dos razones. La primera es empírica: no hay motivo previo para creer que un
sistema de lenguaje y un estudiante de segundo año exhiban la misma dinámica de
negociación, y de hecho hay razones para esperar que no --- un modelo sin costo
de escritura ni pudor por parecer ignorante tiene incentivos comunicativos
distintos a los de una persona. La segunda es conceptual: reducir una modalidad a
estimador de la otra impide preguntarse en qué difieren, que es donde está el
contenido interesante.

Proponemos entonces una formulación simétrica. Sea $\mathcal{P}$ una instancia y
sea $m$ la modalidad. Medimos el mismo perfil de indicadores bajo ambas y
estudiamos el mapa que las vincula. La pregunta central es:

\begin{quote}
\emph{¿En qué medida una misma instancia con partición de información induce la
misma estructura colaborativa cuando el interlocutor es un agente artificial y
cuando es un estudiante universitario, y cómo se descompone la diferencia cuando
existe?}
\end{quote}

Llamamos a esa propiedad \textbf{compatibilidad de modalidad}. No es una
pregunta binaria: una instancia puede ser compatible en la calidad del proceso e
incompatible en la eficiencia, o preservar el ordenamiento relativo entre
problemas sin preservar los niveles absolutos. El marco que sigue está construido
para distinguir esos casos en lugar de promediarlos.

\paragraph{Organización y contribuciones.}
La §\ref{sec:antecedentes} sitúa el trabajo. La §\ref{sec:marco} reformula la
batería de indicadores de modo que sea empíricamente estimable, con dos resultados
propios: la normalización de la interdependencia epistémica, que separa
dificultad de colaboratividad (Observación~\ref{obs:normalizacion}), y la forma
cerrada del balance de carga (Proposición~\ref{prop:ibc}), que elimina la entropía
del marco. La §\ref{sec:compat} construye la noción de compatibilidad, su
descomposición por equiparación lineal (Proposición~\ref{prop:descomp}) y el
límite de resolución del ordenamiento (Proposición~\ref{prop:resolucion}). La
§\ref{sec:diseno} especifica un diseño simétrico con su potencia. Las
§§\ref{sec:amenazas}--\ref{sec:etica} tratan validez y ética, y la
§\ref{sec:implementacion} la instrumentación.

%==============================================================================
\section{Antecedentes}\label{sec:antecedentes}
%==============================================================================

\pendiente{esta sección se redacta sin referencias por decisión de etapa. Las
atribuciones deben incorporarse antes de cualquier circulación externa. Los cuatro
cuerpos de literatura que corresponde acreditar son: (i) el paradigma de perfil
oculto en psicología social, que introduce la estructura de información no
compartida y documenta el fracaso sistemático de los grupos en ponerla en común;
(ii) el marco de dominio de evaluación internacional de resolución colaborativa,
del que se toma la matriz de doce celdas; (iii) los estudios de validación que
compararon compañeros humanos contra agentes guionados en evaluación educativa;
(iv) el trabajo reciente sobre agentes de lenguaje bajo asimetría de información y
sobre el uso de modelos como jueces automáticos.}

Tres observaciones de posicionamiento, independientes de la acreditación
bibliográfica, delimitan el aporte.

En primer lugar, la estructura de información repartida no es nueva; lo nuevo es
el signo con que se la usa. La tradición que la introdujo la empleaba para
documentar un fracaso: los grupos tienden a discutir lo que ya comparten e ignorar
lo que sólo uno sabe. En diseño instruccional se la emplea con el signo opuesto,
para \emph{forzar} la puesta en común. La diferencia importa porque los
indicadores heredados de aquella tradición están calibrados para detectar la
ausencia de intercambio, no su calidad.

En segundo lugar, la sustitución de compañeros humanos por agentes en evaluación
educativa ya se ensayó con agentes guionados por reglas, cuyo comportamiento es
por construcción estable y auditable. Un agente de lenguaje no ofrece ninguna de
las dos garantías: su conducta varía entre ejecuciones y no es inspeccionable. Los
resultados de aquella línea no se transfieren, y es precisamente por eso que la
compatibilidad debe volver a establecerse empíricamente.

En tercer lugar, el marco de partida de este trabajo --- un conjunto de medidores
de interdependencia epistémica y profundidad colaborativa --- fue formulado en
términos probabilísticos y de teoría de la información sin especificar el espacio
de probabilidad ni el estimador. La §\ref{sec:marco} lo reconstruye conservando la
intuición y reemplazando las definiciones no operacionales.

%==============================================================================
\section{Marco de medición}\label{sec:marco}
%==============================================================================

\subsection{Instancias con partición}

\begin{definicion}[Instancia CPS]\label{def:instancia}
Una \emph{instancia con partición de información} es una tupla
$\mathcal{P}=(Q,d_A,d_B,a^\ast,\mathcal{A})$, donde $Q$ es el enunciado público,
$d_A$ y $d_B$ son los datos privados, $\mathcal{A}$ el espacio de respuestas
admisibles y $a^\ast\in\mathcal{A}$ la respuesta correcta, única. Para
$S\subseteq D:=\{d_A,d_B\}$ escribimos $\Info(S):=Q\cup S$.
\end{definicion}

La unicidad de $a^\ast$ no es un detalle: sin ella la corrección no es verificable
por máquina y toda la cadena de medición queda supeditada a juicio humano.

\subsection{Solver de referencia}

El indicador principal del marco de partida se apoya en una cantidad de la forma
$P(\text{solución}\mid \Info)$. Para un problema con respuesta única esa
probabilidad no está definida hasta que se especifica \emph{quién} intenta
resolverlo. Lo hacemos explícito y, con ello, estimable.

\begin{definicion}[Solver de referencia y función de competencia]\label{def:solver}
Un \emph{solver de referencia} es un mapa estocástico
$\solver:\Info\to\mathcal{A}$ especificado por completo: modelo, versión,
temperatura, prompt y presupuesto de razonamiento. Su \emph{función de
competencia} es el juego cooperativo $(D,v)$ dado por
\[
  v(S)\;:=\;\Pr\big[\solver(\Info(S))=a^\ast\big],
  \qquad v:2^{D}\longrightarrow[0,1].
\]
\end{definicion}

\begin{definicion}[Estimador]\label{def:estimador}
Con $N$ ejecuciones independientes sobre $\Info(S)$,
\[
  \hat v(S)=\frac1N\sum_{j=1}^{N}\ind{\solver_j(\Info(S))=a^\ast},
  \qquad N\hat v(S)\sim\mathrm{Bin}\big(N,v(S)\big),
\]
con intervalo de Wilson. La corrección se decide por equivalencia simbólica contra
$a^\ast$; nunca por autorreporte del solver.
\end{definicion}

\begin{observacion}[Relatividad al solver]
$v$ depende de $\solver$, igual que toda medida de dificultad depende de la
población de resolutores sobre la que se calibra. Lo exigible no es eliminar esa
dependencia sino declararla y acotarla, verificando la estabilidad del
\emph{ordenamiento} de instancias entre solvers distintos (§\ref{sec:amenazas}).
Cuando la modalidad involucra estudiantes, la función de competencia relevante es
la de una población humana, $v^{H}$, que se estima en un piloto separado
(§\ref{sec:piloto}).
\end{observacion}

\begin{supuesto}[Monotonía informacional]\label{sup:monotonia}
$v(\emptyset)\le v(\{d_i\})\le v(D)$ para $i\in\{A,B\}$.
\end{supuesto}

El supuesto es normativo: información adicional no debería perjudicar. Los
sistemas reales pueden violarlo por distracción. Las violaciones se reportan, no
se descartan: son diagnóstico sobre la instancia o sobre el solver.

\subsection{Perfil de diseño}

Los indicadores de esta subsección dependen sólo de $\mathcal{P}$ y de $\solver$,
no del transcript. Son, por construcción, independientes de la modalidad, y por
eso constituyen la base común sobre la cual la comparación de §\ref{sec:compat}
tiene sentido.

\begin{definicion}[Ganancia alcanzable]
$\Delta:=v(D)-v(\emptyset)$. La instancia es \emph{no degenerada} si $\Delta>0$.
\end{definicion}

\begin{definicion}[Interdependencia epistémica de diseño]\label{def:ie0}
\[
  IE_0\;:=\;1-\frac{\max_{i}\big[v(\{d_i\})-v(\emptyset)\big]}{\Delta}.
\]
\end{definicion}

\begin{proposicion}\label{prop:ie}
Bajo el Supuesto~\ref{sup:monotonia} y no degeneración:
(i) $IE_0\in[0,1]$;
(ii) $IE_0=1$ si y sólo si $v(\{d_A\})=v(\{d_B\})=v(\emptyset)$;
(iii) $IE_0=0$ si y sólo si $\max_i v(\{d_i\})=v(D)$.
\end{proposicion}

\begin{proof}
Por monotonía, $0\le v(\{d_i\})-v(\emptyset)\le\Delta$ para cada $i$; tomando
máximo y dividiendo por $\Delta>0$, el sustraendo pertenece a $[0,1]$, lo que da
(i). Los incisos (ii) y (iii) son los casos de igualdad en cada extremo.
\end{proof}

\begin{observacion}[La normalización no es cosmética]\label{obs:normalizacion}
El indicador original, $1-\max_i P(\text{sol}\mid I_i)$, confunde
\emph{dificultad} con \emph{interdependencia}: una instancia que nadie resuelve ni
con ambos datos ($v(D)=v(\{d_i\})=0$) alcanza el valor máximo sin ser
colaborativa, sólo por ser imposible. Normalizar por la ganancia alcanzable separa
los dos constructos y deja la dificultad como indicador propio, $1-v(D)$, en lugar
de contaminar al principal.
\end{observacion}

\paragraph{Balance de carga.}
Una instancia puede exigir que ambos hablen y ser colaborativamente superficial si
uno aporta casi todo el peso deductivo y el otro un dato trivial. El marco de
partida captura esto con entropías condicionales, lo que introduce dos problemas:
sobre un espacio de respuestas continuo la entropía diferencial depende de la
escala y puede ser negativa, de modo que el índice no queda acotado; y no es
estimable con el mismo experimento que $IE_0$. Ambos desaparecen al usar el valor
de Shapley del juego $(D,v)$.

\begin{definicion}[Contribución marginal]
Para $|D|=2$,
\[
  \varphi_A=\tfrac12\big[v(\{d_A\})-v(\emptyset)\big]+\tfrac12\big[v(D)-v(\{d_B\})\big],
\]
y simétricamente $\varphi_B$.
\end{definicion}

\begin{proposicion}\label{prop:ibc}
$\;$(i) $\varphi_A+\varphi_B=\Delta$; \quad
(ii) $\varphi_A-\varphi_B=v(\{d_A\})-v(\{d_B\})$.
\end{proposicion}

\begin{proof}
Sumando ambas expresiones,
\[
  \varphi_A+\varphi_B
  =\tfrac12\big[v(\{d_A\})+v(\{d_B\})-2v(\emptyset)\big]
  +\tfrac12\big[2v(D)-v(\{d_A\})-v(\{d_B\})\big]
  =v(D)-v(\emptyset)=\Delta,
\]
que es (i). Restando,
\[
  \varphi_A-\varphi_B
  =\tfrac12\big[v(\{d_A\})-v(\{d_B\})\big]+\tfrac12\big[v(\{d_A\})-v(\{d_B\})\big]
  =v(\{d_A\})-v(\{d_B\}).\qedhere
\]
\end{proof}

\begin{definicion}[Índice de balance de carga informativa]
\[
  IBC\;:=\;1-\frac{|\varphi_A-\varphi_B|}{\Delta}
        \;=\;1-\frac{\big|v(\{d_A\})-v(\{d_B\})\big|}{v(D)-v(\emptyset)}.
\]
\end{definicion}

\begin{corolario}
Bajo el Supuesto~\ref{sup:monotonia} y no degeneración, $IBC\in[0,1]$; vale $1$ si
y sólo si ambos datos producen idéntica ganancia individual, y $0$ si y sólo si
uno concentra toda la ganancia alcanzable.
\end{corolario}

\begin{proof}
Sin pérdida de generalidad $v(\{d_A\})\ge v(\{d_B\})$. Por monotonía
$v(\{d_A\})\le v(D)$ y $v(\{d_B\})\ge v(\emptyset)$, de donde
$0\le v(\{d_A\})-v(\{d_B\})\le v(D)-v(\emptyset)=\Delta$; dividiendo por $\Delta$
se obtiene el rango, y los casos de igualdad dan las caracterizaciones.
\end{proof}

El inciso (ii) de la Proposición~\ref{prop:ibc} es el resultado operativamente
útil: el $IBC$ se calcula con las mismas cuatro proporciones ya estimadas para
$IE_0$, sin experimento adicional y sin entropías.

\begin{observacion}[Más de dos participantes]
Con $k$ agentes, $\varphi_i$ recupera la forma general de Shapley y el índice se
generaliza como $IBC_k=1-(\max_i\varphi_i-\min_i\varphi_i)/\Delta$, o mediante el
coeficiente de Gini de $(\varphi_1,\dots,\varphi_k)$ si interesa penalizar
desbalances distribuidos y no sólo el rango.
\end{observacion}

\begin{proposicion}[Varianza por método delta]\label{prop:delta}
Sea $g(v)=(v_A-v_B)/(v_D-v_\emptyset)$ con estimaciones independientes de tamaño
$N$ por celda. Entonces
\[
  \Var(\hat g)\;\approx\;
  \frac{v_A(1-v_A)+v_B(1-v_B)}{N\Delta^{2}}
  +\frac{(v_A-v_B)^{2}\big[v_D(1-v_D)+v_\emptyset(1-v_\emptyset)\big]}{N\Delta^{4}}.
\]
\end{proposicion}

\begin{proof}
Las derivadas parciales de $g$ son $1/\Delta$, $-1/\Delta$,
$-(v_A-v_B)/\Delta^{2}$ y $(v_A-v_B)/\Delta^{2}$ respecto de $v_A$, $v_B$, $v_D$ y
$v_\emptyset$. El método delta con $\Var(\hat v_S)=v_S(1-v_S)/N$ e independencia
entre celdas da la expresión.
\end{proof}

El término en $\Delta^{-4}$ advierte que la estimación se degrada cuando la
ganancia alcanzable es pequeña; con $\hat\Delta$ chico se usan intervalos
bootstrap BCa en lugar de la aproximación normal.

\paragraph{Tamaño muestral por celda.} Verificar la insuficiencia individual
($v(\{d_i\})\le\delta_0$) es un problema de cota superior con cero eventos. Si
ninguna de $N$ ejecuciones acierta, la cota unilateral exacta al $95\%$ es
$1-0.05^{1/N}$, esto es $0.095$ para $N=30$ y $0.049$ para $N=60$. Fijamos $N=30$
por celda, elevado a $60$ en instancias limítrofes.

\paragraph{Estructura de la solución.}
La resolución canónica se representa como un DAG $G=(V,E)$ de pasos deductivos;
$\Cl(S)\subseteq V$ denota los nodos derivables a partir de $\Info(S)$. Entonces
\[
  MEE:=\frac{\big|V\setminus(\Cl(\{d_A\})\cup\Cl(\{d_B\}))\big|}
            {\big|V\setminus\Cl(\emptyset)\big|}\in[0,1]
\]
mide la fracción de pasos no triviales que sólo existen como producto de la
deducción conjunta. Asignando a cada nodo el lado $\lambda(\nu)\in\{A,B,AB\}$ del
dato privado mínimo que su derivación requiere,
\[
  t_{\min}:=\max_{\pi\in\mathrm{caminos}(G)}
  \#\big\{j:\lambda(\pi_j)\neq\lambda(\pi_{j+1}),\;
  \lambda(\pi_j),\lambda(\pi_{j+1})\in\{A,B\}\big\}
\]
cuenta el máximo de alternancias de lado a lo largo de un camino dirigido, y es
una cota inferior estructural sobre el número de intercambios necesarios: cada
alternancia en la cadena de dependencias exige al menos una transmisión.

\begin{definicion}[Perfil de diseño]\label{def:phi}
\[
  \Phi(\mathcal{P})=\big(IE_0,\;IBC,\;MEE,\;t_{\min},\;1-v(D)\big).
\]
\end{definicion}

\subsection{Perfil de episodio}

Los indicadores que siguen se computan sobre un transcript y por lo tanto
\emph{sí} dependen de la modalidad. Son el objeto de la comparación.

Sea $m_{1:t}$ el transcript hasta la ronda $t$ e $I_{i,t}=Q\cup\{d_i\}\cup m_{1:t}$.
Con $v^{(i)}_t=\Pr[\solver(I_{i,t})=a^\ast]$,
\[
  IE_t=\clip_{[0,1]}\!\left(1-\frac{\max_i v^{(i)}_t-v(\emptyset)}{\Delta}\right),
  \qquad
  t_{\mathrm{crit}}=\max\{t\le T: IE_t\ge \tfrac12\},
\]
con la convención $\max\emptyset=0$. Una instancia cuya partición se disuelve en el
primer intercambio tiene $t_{\mathrm{crit}}=0$.

\begin{definicion}[Calidad del proceso]
Sobre las doce celdas del marco de dominio, con $q_k\in\{0,1,2,3\}$ en escala
ordinal desde ausencia de interacción hasta emergencia de entendimiento conjunto,
$CQI:=\tfrac1{36}\sum_{k=1}^{12}q_k\in[0,1]$.
\end{definicion}

\begin{definicion}[Eficiencia y divulgación prematura]
$E:=t_{\min}/t_{\mathrm{obs}}\in(0,1]$, con $t_{\mathrm{obs}}$ las rondas usadas
hasta la primera respuesta correcta. Con
$T_i:=\min\{t: d_i \text{ es recuperable en su totalidad de } m_{1:t}\}$, el
índice de divulgación prematura es $PD_i:=\ind{T_i=1}$.
\end{definicion}

Sostenemos que la divulgación prematura debe medirse y no prohibirse por prompt.
Una restricción explícita sustituye el fenómeno de interés --- el comportamiento
comunicativo espontáneo del participante --- por otro --- su obediencia a la
restricción --- y con ello destruye la comparabilidad entre modalidades, que es
justamente el objeto del trabajo. La restricción se incorpora como condición
experimental cruzada.

\begin{definicion}[Perfil de episodio]\label{def:psi}
\[
  \Psi(\omega)=\big(\underbrace{t_{\mathrm{crit}},\;CQI,\;E}_{\text{continuos}},\;
  \underbrace{\ind{\text{correcto}},\;PD}_{\text{binarios}}\big)\in\Real^{J},\quad J=5.
\]
\end{definicion}

\subsection{Requisitos de medición}

\begin{requisito}[Corrección verificada]\label{req:correccion}
$\ind{\text{correcto}}$ se determina por equivalencia simbólica de la respuesta
final contra $a^\ast$, mediante un sistema de álgebra computacional con
normalización previa. No se acepta en ningún caso la autodeclaración de resolución
por parte del agente o del participante.
\end{requisito}

El Requisito~\ref{req:correccion} no es pedantería. Un marcador de resolución
autodeclarado es falsificable por el propio sujeto medido y convierte la variable
dependiente principal en una medida de disposición a declararse exitoso.

\begin{requisito}[Codificación del proceso]\label{req:juez}
$CQI$ se obtiene por codificación automática mediante modelos jueces, con:
(a) al menos dos jueces de familias distintas, ninguno de la familia que generó
los transcripts;
(b) ceguera respecto de la modalidad, la condición y la identidad de la instancia;
(c) orden de presentación aleatorizado;
(d) acuerdo juez--humano reportado como $\alpha$ de Krippendorff ordinal sobre una
submuestra estratificada del $20\%$ codificada por dos anotadores entrenados,
con los umbrales convencionales $\alpha\ge0.80$ para conclusiones firmes y
$0.667\le\alpha<0.80$ para conclusiones tentativas.
\end{requisito}

La ceguera respecto de la modalidad en (b) es crítica y no trivial de implementar:
los transcripts humanos y artificiales difieren en superficie (ortografía,
longitud, muletillas) de un modo que puede delatar la condición al juez. El
protocolo incluye una normalización estilística previa y una prueba de
desenmascaramiento en la que se pide al juez identificar la modalidad; si acierta
por encima del azar, la ceguera falló y el resultado debe reportarse con esa
salvedad.

\subsection{Sobre los índices compuestos}\label{sec:compuesto}

El marco de partida propone un rendimiento colaborativo de la forma
$CY=0.35\,CDI+0.45\,\ind{\text{correcto}}+0.20\,\kappa$. Señalamos dos defectos.

\begin{observacion}[Varianza intra-instancia nula]
$CDI$ es, por definición, prescriptivo: función de $\mathcal{P}$ y no del episodio
$\omega$. Luego
$\Var_\omega(CY\mid\mathcal{P})=\Var_\omega(0.45\,\ind{\text{correcto}}+0.20\kappa\mid\mathcal{P})$:
el término en $CDI$ no aporta varianza dentro de una instancia y opera sólo como
desplazamiento entre instancias. Un índice de rendimiento a nivel de episodio no
puede contener un término constante en el episodio.
\end{observacion}

\begin{observacion}[Escalarización con pérdida]
Si calidad de proceso y exactitud son ortogonales --- como el propio marco de
partida afirma --- una suma ponderada de pesos fijos las colapsa en un escalar no
identificable: infinitos pares producen el mismo valor y la ortogonalidad
garantiza que la pérdida no es recuperable. El parámetro $\kappa$, además, no está
definido.
\end{observacion}

Por eso reportamos $\Phi$ y $\Psi$ como vectores. Cualquier compuesto es
secundario y va acompañado de análisis de sensibilidad sobre sus pesos. Reportamos
además la matriz de correlaciones de $\Phi$ y sus componentes principales: si la
primera componente explica más del $80\%$ de la varianza, el marco tiene
indicadores redundantes y debe reducirse.

%==============================================================================
\section{Compatibilidad entre modalidades}\label{sec:compat}
%==============================================================================

\subsection{La modalidad como factor}

\begin{definicion}[Modalidad]
$m\in\mathcal{M}=\{\modAA,\modHA\}$, donde \modAA{} designa el diálogo entre dos
agentes artificiales y \modHA{} el diálogo entre un agente artificial y un
estudiante. Un episodio bajo modalidad $m$ sobre la instancia $\mathcal{P}$ es una
realización $\omega\sim\mathcal{D}_m(\mathcal{P})$, y
\[
  \mu_{m,j}(\mathcal{P}):=\Esp_{\omega\sim\mathcal{D}_m(\mathcal{P})}\big[\Psi_j(\omega)\big],
  \qquad j=1,\dots,J.
\]
\end{definicion}

Nótese que $\Phi(\mathcal{P})$ no lleva subíndice de modalidad: es la base común
que hace comparables a $\mu_{\modAA}$ y $\mu_{\modHA}$. Sin un perfil de diseño
independiente de la modalidad, cualquier diferencia observada sería atribuible a
que se trata, en rigor, de dos instancias distintas.

\subsection{Divergencia estandarizada}

\begin{definicion}[Divergencia por componente]\label{def:divergencia}
Para componentes continuos, sea $\bar\sigma_j$ la desviación estándar
intra-celda agrupada sobre todo el banco y ambas modalidades. La \emph{divergencia
estandarizada} de la instancia $\mathcal{P}$ en el componente $j$ es
\[
  d_j(\mathcal{P})\;:=\;\frac{\mu_{\modAA,j}(\mathcal{P})-\mu_{\modHA,j}(\mathcal{P})}{\bar\sigma_j}.
\]
Para componentes binarios se usa la diferencia de proporciones en la escala de
probabilidad, $\pi_j(\mathcal{P}):=\mu_{\modAA,j}(\mathcal{P})-\mu_{\modHA,j}(\mathcal{P})$,
con su propio margen de equivalencia.
\end{definicion}

Agrupar $\bar\sigma_j$ sobre el banco, y no dentro de cada instancia, es
deliberado: vuelve $d_j$ comparable entre instancias y evita que una instancia con
varianza interna pequeña produzca divergencias estandarizadas artificialmente
grandes. La distinción entre componentes continuos y binarios tampoco es
cosmética: una diferencia estandarizada sobre una variable de Bernoulli confunde
el tamaño del efecto con la tasa base.

\subsection{Compatibilidad absoluta y relativa}

Dos instancias pueden diferir en niveles y coincidir en ordenamiento. Como son
propiedades distintas y de utilidad distinta, las separamos.

\begin{definicion}[Compatibilidad absoluta]\label{def:absoluta}
Dado un margen $\delta>0$, la instancia $\mathcal{P}$ es
\emph{$(j,\delta)$-compatible} si $|d_j(\mathcal{P})|\le\delta$. El
\emph{índice de compatibilidad} de la instancia es la fracción de componentes
declarados equivalentes,
\[
  C_\delta(\mathcal{P})=\frac1J\sum_{j=1}^{J}\ind{\mathcal{P}\text{ es }(j,\delta)\text{-compatible}}\in[0,1],
\]
y su \emph{divergencia máxima} es
$\mathrm{Div}(\mathcal{P})=\max_j|d_j(\mathcal{P})|$.
\end{definicion}

\begin{definicion}[Compatibilidad relativa]\label{def:relativa}
El banco es \emph{relativamente compatible} en el componente $j$ si el coeficiente
de correlación de rangos de Spearman entre
$\mu_{\modAA,j}(\cdot)$ y $\mu_{\modHA,j}(\cdot)$, calculado sobre las instancias
del banco, supera un umbral $\rho_0$ fijado de antemano.
\end{definicion}

La compatibilidad relativa es la propiedad que interesa a quien construye un banco
de problemas: basta con que las instancias se ordenen igual bajo ambas modalidades
para poder seleccionarlas con la modalidad barata. La compatibilidad absoluta es
la que interesa a quien quiere interpretar un valor. La primera no implica la
segunda, y la segunda implica la primera sólo hasta la resolución que cuantifica
el siguiente resultado.

\begin{proposicion}[Resolución del ordenamiento]\label{prop:resolucion}
Supóngase $|d_j(\mathcal{P})|\le\delta$ para toda instancia del banco. Entonces,
para dos instancias $\mathcal{P},\mathcal{P}'$, el signo de
$\mu_{\modAA,j}(\mathcal{P})-\mu_{\modAA,j}(\mathcal{P}')$ coincide con el de
$\mu_{\modHA,j}(\mathcal{P})-\mu_{\modHA,j}(\mathcal{P}')$ siempre que
\[
  \frac{\big|\mu_{\modHA,j}(\mathcal{P})-\mu_{\modHA,j}(\mathcal{P}')\big|}{\bar\sigma_j}\;>\;2\delta .
\]
\end{proposicion}

\begin{proof}
Escribiendo $\mu_{\modAA,j}=\mu_{\modHA,j}+\bar\sigma_j d_j$,
\[
  \mu_{\modAA,j}(\mathcal{P})-\mu_{\modAA,j}(\mathcal{P}')
  =\big[\mu_{\modHA,j}(\mathcal{P})-\mu_{\modHA,j}(\mathcal{P}')\big]
  +\bar\sigma_j\big[d_j(\mathcal{P})-d_j(\mathcal{P}')\big].
\]
Por hipótesis $|d_j(\mathcal{P})-d_j(\mathcal{P}')|\le 2\delta$, de modo que el
segundo sumando está acotado en valor absoluto por $2\delta\bar\sigma_j$. Si el
primero excede esa cota, domina y fija el signo.
\end{proof}

\begin{observacion}[Cuántos estratos son distinguibles]\label{obs:estratos}
Si la dispersión entre instancias de $\mu_{\modHA,j}$ es $s_j$ en unidades de
$\bar\sigma_j$, la Proposición~\ref{prop:resolucion} implica que el número de
estratos que la modalidad barata puede distinguir con confianza es del orden de
$s_j/(2\delta)$. Con una dispersión típica de $1.5\,\bar\sigma_j$ y un margen
$\delta=0.5$, eso da apenas $1.5$: el cribado permite separar dos grupos, no
producir un ranking fino. Es una limitación estructural del procedimiento y debe
declararse como tal en lugar de descubrirse al interpretar los resultados.
\end{observacion}

\subsection{Descomposición por equiparación lineal}

Cuando hay incompatibilidad conviene saber de qué tipo es. Postulamos, para cada
componente $j$, el modelo de equiparación sobre las instancias del banco
\begin{equation}\label{eq:equiparacion}
  \mu_{\modAA,j}(\mathcal{P})=a_j+b_j\,\mu_{\modHA,j}(\mathcal{P})+e_j(\mathcal{P}),
  \qquad \Esp[e_j]=0,\quad e_j\perp\mu_{\modHA,j}.
\end{equation}

\begin{proposicion}[Descomposición de la divergencia esperada]\label{prop:descomp}
Bajo \eqref{eq:equiparacion}, escribiendo $X=\mu_{\modHA,j}(\mathcal{P})$ con
$\mathcal{P}$ recorriendo el banco,
\[
  \Esp\big[(\mu_{\modAA,j}-\mu_{\modHA,j})^{2}\big]
  =\underbrace{\big(a_j+(b_j-1)\Esp[X]\big)^{2}}_{\text{sesgo sistemático}}
  +\underbrace{(b_j-1)^{2}\Var(X)}_{\text{distorsión de escala}}
  +\underbrace{\Var(e_j)}_{\text{idiosincrásico}} .
\]
\end{proposicion}

\begin{proof}
Sea $Y=\mu_{\modAA,j}$. De \eqref{eq:equiparacion},
$Y-X=a_j+(b_j-1)X+e_j$, luego
$\Esp[Y-X]=a_j+(b_j-1)\Esp[X]$ y, por independencia de $e_j$ y $X$,
$\Var(Y-X)=(b_j-1)^{2}\Var(X)+\Var(e_j)$. La identidad
$\Esp[(Y-X)^2]=\Var(Y-X)+(\Esp[Y-X])^{2}$ concluye.
\end{proof}

La lectura de los tres términos es directa y sugiere remedios distintos. Un
\emph{sesgo sistemático} significa que una modalidad es uniformemente más (o
menos) colaborativa que la otra; es corregible por desplazamiento y no invalida el
uso comparativo. Una \emph{distorsión de escala} con $b_j<1$ indica compresión de
rango: la modalidad artificial discrimina menos entre instancias de lo que
discriminan los estudiantes, lo que degrada el ordenamiento sin desplazarlo. El
término \emph{idiosincrásico} es el único irreducible: mide en qué medida la
divergencia depende de la instancia particular y no de una transformación global.

\begin{observacion}[Techo de confiabilidad]
$\Var(e_j)$ no puede caer por debajo de la varianza de error de medición de ambos
lados. Reportamos el cociente entre $\Var(e_j)$ y esa varianza de error estimada;
si es cercano a uno, toda la incompatibilidad idiosincrásica observada es ruido y
no señal.
\end{observacion}

\subsection{Inferencia: equivalencia, no ausencia de diferencia}

\begin{observacion}[El error inferencial que este diseño evita]\label{obs:tost}
No rechazar la hipótesis nula de diferencia \emph{no} es evidencia de
equivalencia: un estudio con poca potencia produce sistemáticamente ese resultado.
Afirmar compatibilidad exige invertir la carga de la prueba.
\end{observacion}

\begin{definicion}[Prueba de equivalencia]\label{def:tost}
Fijado el margen $\delta$, se declara $(j,\delta)$-compatibilidad si el intervalo
de confianza al $(1-2\alpha)$ para $d_j(\mathcal{P})$ está contenido en
$[-\delta,\delta]$; equivalentemente, si ambas hipótesis unilaterales
$H_{01}: d_j\le-\delta$ y $H_{02}: d_j\ge\delta$ se rechazan al nivel $\alpha$.
\end{definicion}

\begin{proposicion}[Tamaño muestral para equivalencia]\label{prop:potencia}
Para dos grupos independientes de igual tamaño, con divergencia verdadera nula,
nivel $\alpha$ y potencia $1-\beta$, el número de episodios por modalidad es
aproximadamente
\[
  n\;\approx\;\frac{2\big(z_{1-\alpha}+z_{1-\beta/2}\big)^{2}}{\delta^{2}} .
\]
\end{proposicion}

Con $\alpha=0.05$ y potencia $0.80$ la fórmula da los valores de la
Tabla~\ref{tab:potencia}, confirmados por simulación. Son números grandes, y esa
es exactamente la información que el diseño necesita antes de recolectar: la
equivalencia a nivel de instancia individual es cara, y por lo tanto sólo puede
sostenerse sobre un subconjunto reducido de instancias.

\begin{table}[t]
\centering
\begin{tabular}{@{}cccccc@{}}
\toprule
Margen $\delta$ & $0.3$ & $0.4$ & $0.5$ & $0.6$ & $0.8$ \\
\midrule
Episodios por modalidad & $191$ & $108$ & $69$ & $48$ & $27$ \\
\bottomrule
\end{tabular}
\caption{Episodios necesarios por modalidad para declarar equivalencia a nivel de
instancia ($\alpha=0.05$, potencia $0.80$, divergencia verdadera nula).}
\label{tab:potencia}
\end{table}

%==============================================================================
\section{Diseño experimental}\label{sec:diseno}
%==============================================================================

El diseño es deliberadamente simétrico en su instrumento --- ambas modalidades se
miden con el mismo $\Psi$ y el mismo protocolo de codificación --- y asimétrico en
su volumen, porque el costo por episodio difiere en tres órdenes de magnitud.

\subsection{Materiales}

$50$ instancias distribuidas sobre seis áreas de matemática universitaria:
álgebra lineal, introducción al álgebra, introducción al cálculo, cálculo
diferencial e integral, cálculo en varias variables y ecuaciones diferenciales
ordinarias. Cada una con enunciado público, partición $(d_A,d_B)$, respuesta
canónica estructurada $a^\ast$ y DAG de solución anotado, del que se derivan
$MEE$ y $t_{\min}$.

\subsection{Estudio 1: modalidad agente--agente}

Diseño factorial intra-instancia $2\times2$: familia de modelo $\{M_1,M_2\}$
cruzada con restricción de divulgación $\{$libre, restringida$\}$, con $R=8$
réplicas por celda. Total: $50\times4\times8=1600$ episodios.

\subsection{Estudio 2: modalidad agente--estudiante}\label{sec:estudio2}

Dos brazos, porque las dos formas de compatibilidad tienen requerimientos
muestrales incompatibles entre sí.

\paragraph{Brazo extensivo (equiparación).} $25$ instancias, submuestreadas de
forma estratificada por terciles de $IE_0$ y $t_{\min}$ para cubrir el rango del
banco y no sólo su extremo favorable. Aproximadamente $8$ episodios por instancia,
esto es $200$ episodios. La unidad de análisis es la instancia, $n=25$, y el
objetivo es estimar $a_j$, $b_j$ y $\Var(e_j)$ de \eqref{eq:equiparacion} y la
correlación de rangos de la Definición~\ref{def:relativa}.

\paragraph{Brazo intensivo (equivalencia).} Un subconjunto de $6$ instancias con
$n=48$ episodios cada una, dimensionado según la Tabla~\ref{tab:potencia} para un
margen $\delta=0.6$: $288$ episodios. Es el único brazo que autoriza afirmaciones
de equivalencia a nivel de instancia, y sólo con ese margen.

\paragraph{Participantes.} Estudiantes de grado con la materia correspondiente
aprobada o en curso, cada uno en $3$ episodios, con orden contrabalanceado por
cuadrado latino. El total de $488$ episodios requiere del orden de $163$
participaciones.

\paragraph{Contrabalanceo de rol.} En cada instancia, la mitad de los episodios
asigna $d_A$ al estudiante y la otra mitad $d_B$. Sin esto, la comparación entre
modalidades confunde el efecto de la modalidad con el del rol, dado que en
\modAA{} ambos roles son artificiales mientras que en \modHA{} el humano ocuparía
siempre el mismo. Es un requisito, no un refinamiento.

\subsection{Piloto de insuficiencia individual humana}\label{sec:piloto}

Toda la construcción presupone que ningún dato alcanza por separado. Esa premisa
se verifica para el solver artificial por la Definición~\ref{def:estimador}, pero
no para estudiantes. Un piloto independiente con $12$ instancias y estudiantes que
reciben un solo dato estima $v^{H}(\{d_i\})$ y verifica que la insuficiencia
individual valga también para la población humana. Si no valiera, la instancia no
es colaborativa para estudiantes por más que lo sea para el solver, y la
comparación de modalidades pierde sentido para esa instancia.

\subsection{Hipótesis}

\begin{description}
\item[H1 (compatibilidad relativa).] Para $CQI$ y $t_{\mathrm{crit}}$, la
  correlación de rangos entre modalidades sobre las $25$ instancias del brazo
  extensivo es positiva y supera $\rho_0=0.5$.
\item[H2 (estructura de la divergencia).] En la descomposición de la
  Proposición~\ref{prop:descomp}, el término idiosincrásico no domina: representa
  menos de la mitad de la divergencia cuadrática esperada, de modo que la
  incompatibilidad es mayoritariamente corregible por transformación afín.
\item[H3 (compatibilidad absoluta acotada).] En el brazo intensivo, al menos tres
  de los cinco componentes de $\Psi$ resultan $(j,0.6)$-compatibles por la
  Definición~\ref{def:tost}.
\item[H4 (divergencia en divulgación).] La tasa de divulgación prematura
  $\Esp[PD]$ es estrictamente mayor en \modAA{} que en \modHA{}. Hipótesis
  direccional: predice que los agentes colapsan la partición antes que los
  estudiantes.
\item[H5 (compresión de rango).] $b_j<1$ para $CQI$: la modalidad artificial
  discrimina entre instancias menos que la humana.
\end{description}

H4 y H5 son las hipótesis que pueden falsar el programa, y están formuladas para
que su confirmación sea informativa incluso si H1 falla. Si los agentes vuelcan su
dato de entrada mientras los estudiantes negocian, las dos modalidades no miden lo
mismo, y saber \emph{eso} --- con un tamaño de efecto y no como impresión --- es el
resultado que ordena el uso de agentes en investigación educativa.

\subsection{Análisis}

El modelo principal del Estudio 2 es de efectos aleatorios cruzados,
\[
  \Psi_{j,sp}=\beta_{0}+\boldsymbol{\beta}^{\!\top}\Phi_p+u_s+w_p+\varepsilon_{sp},
  \qquad u_s\sim\mathcal{N}(0,\sigma_s^{2}),\;\; w_p\sim\mathcal{N}(0,\sigma_p^{2}),
\]
con estudiantes ($s$) e instancias ($p$) como factores aleatorios cruzados. Los
contrastes entre modalidades se estiman sobre las medias por instancia, ajustadas
por el modelo.

\paragraph{Confiabilidad de las medias por instancia.} Con $k$ episodios por
instancia, la confiabilidad de la media sigue la relación de Spearman--Brown
$\rho_k=k\,\mathrm{ICC}/[1+(k-1)\mathrm{ICC}]$. En el brazo extensivo, con
$k=8$ e $\mathrm{ICC}=0.30$, resulta $\rho_k\approx0.77$, suficiente para
correlacionar medias pero no para interpretar valores individuales.

\paragraph{Corrección por atenuación.} Ambos términos de H1 son medidas con error,
de modo que la correlación observada subestima la verdadera. Reportamos también
$\rho_{\text{corr}}=\rho_{\text{obs}}/\sqrt{\rho_{\modAA}\rho_{\modHA}}$, con
$\rho_{\bullet}$ las confiabilidades estimadas de cada lado, y explicitamos ambas
cifras.

%==============================================================================
\section{Amenazas a la validez}\label{sec:amenazas}
%==============================================================================

\paragraph{Contaminación del corpus de entrenamiento.} Si una instancia aparece en
los datos de preentrenamiento, $v(\{d_i\})$ se infla y $IE_0$ se subestima.
Mitigaciones: instancias de autoría propia; prueba canario que perturba constantes
numéricas y verifica que $\hat v$ se mueva como predice la estructura; y reporte
de $\hat v(\emptyset)$, cuyo exceso sobre el azar es señal directa de
contaminación o de un enunciado que filtra la respuesta.

\paragraph{Relatividad al solver.} Todo $\Phi$ depende de $\solver$. Se replica
con al menos dos solvers de familias distintas y se reporta la concordancia de
ordenamientos por $\tau$ de Kendall. Si es baja, los indicadores describen al
modelo y no a la instancia, y el programa debe reformularse.

\paragraph{Confusión modalidad $\times$ rol.} Tratada por contrabalanceo en
§\ref{sec:estudio2}.

\paragraph{Fallo de ceguera del juez.} Tratado por normalización estilística y
prueba de desenmascaramiento en el Requisito~\ref{req:juez}.

\paragraph{Sesgo de verbosidad.} Se reporta la correlación entre longitud del
transcript y $CQI$ asignado. Dado que las modalidades difieren sistemáticamente en
longitud, un sesgo de verbosidad no controlado produciría por sí solo una
divergencia espuria entre modalidades.

\paragraph{Efectos de práctica.} Contrabalanceo por cuadrado latino del orden de
las tres instancias por participante.

\paragraph{Alcance de las afirmaciones.} La Observación~\ref{obs:estratos} y la
Tabla~\ref{tab:potencia} acotan de antemano qué se puede concluir: compatibilidad
relativa sobre $25$ instancias, y equivalencia absoluta sólo sobre $6$ instancias
con margen $0.6$. Ninguna conclusión se extenderá más allá.

%==============================================================================
\section{Consideraciones éticas}\label{sec:etica}
%==============================================================================

Consentimiento informado previo, con descripción del uso de los transcripts.
Participación desacoplada de toda calificación académica: ni el desempeño ni la
negativa a participar tienen consecuencia curricular. Seudonimización en el momento
de la captura, con la tabla de correspondencia almacenada por separado y con
acceso restringido. Derecho de retiro con borrado efectivo de los episodios
asociados. Plazo de retención y condiciones de acceso declarados de antemano.
Sometimiento al comité de ética institucional antes de iniciar el Estudio 2. Se
informa a los participantes que su interlocutor es un sistema artificial; el
engaño sobre ese punto no forma parte del diseño.

%==============================================================================
\section{Instrumentación}\label{sec:implementacion}
%==============================================================================

El protocolo se instrumenta sobre una plataforma web propia que ya soporta las dos
modalidades y persiste transcripts completos con rol y orden. La
Tabla~\ref{tab:brecha} contrasta lo que el diseño exige contra lo disponible.

\begin{table}[t]
\centering
\small
\begin{tabular}{@{}lll@{}}
\toprule
\textbf{Constructo} & \textbf{Requiere} & \textbf{Estado} \\
\midrule
Transcripts por rol y orden          & tabla de mensajes                     & disponible \\
Modalidades \modAA{} y \modHA{}      & dos modos de sesión                   & disponible \\
Partición $d_A/d_B$                  & campos privados por instancia         & disponible \\
$\ind{\text{correcto}}$ (Req.~\ref{req:correccion}) & $a^\ast$ y verificador simbólico & pendiente \\
$v(S)$, $IE_t$, $IBC$                & arnés de solver de referencia         & pendiente \\
$MEE$, $t_{\min}$                    & DAG de solución anotado               & pendiente \\
$CQI$ (Req.~\ref{req:juez})          & jueces, ceguera, desenmascaramiento   & pendiente \\
Contrabalanceo de rol                & asignación de $d_A$/$d_B$ al humano   & pendiente \\
Condiciones experimentales           & aleatorización registrada             & pendiente \\
$E$ y costos                         & tokens y latencia por turno           & pendiente \\
Reproducibilidad                     & versionado de instancias y prompts    & pendiente \\
Análisis                             & exportación en formato tabular        & pendiente \\
Ética                                & consentimiento y seudonimización      & pendiente \\
\bottomrule
\end{tabular}
\caption{Correspondencia entre constructos del marco y capacidades de la
plataforma.}
\label{tab:brecha}
\end{table}

La fila del contrabalanceo de rol merece una nota: la implementación actual asigna
siempre $d_A$ al participante humano, lo que hace estructuralmente imposible el
control descrito en §\ref{sec:estudio2}. Es un ejemplo de cómo una decisión
razonable en un producto se vuelve un defecto de diseño en un instrumento.

%==============================================================================
\section{Discusión}
%==============================================================================

La formulación simétrica tiene una ventaja que la formulación como cribado no
tenía: ningún resultado es un fracaso. Si las modalidades resultan compatibles en
sentido relativo, se dispone de un procedimiento barato para ordenar bancos de
problemas, con el alcance acotado por la Observación~\ref{obs:estratos}. Si
resultan incompatibles, la descomposición de la Proposición~\ref{prop:descomp}
indica de qué tipo es la incompatibilidad y, en dos de los tres casos, cómo
corregirla. Y si el término idiosincrásico domina, se obtiene una medición del
grado en que el comportamiento colaborativo de un sistema de lenguaje difiere del
de un estudiante --- expresada en unidades interpretables y sobre dimensiones
identificadas, no como impresión cualitativa.

Queda fuera del alcance de este diseño una pregunta que lo sucede naturalmente: si
las instancias que resultan compatibles comparten alguna propiedad estructural
reconocible a priori sobre el DAG de solución. De existir, permitiría diseñar por
construcción problemas compatibles en lugar de descubrirlos por medición.

\vspace{1em}
\noindent\pendiente{secciones a incorporar antes de circular: antecedentes con
atribuciones (§\ref{sec:antecedentes}), preinscripción del protocolo, y
disponibilidad de datos y código.}

\end{document}
